% Generated draft: requires mathematical and source review.
\documentclass[11pt]{article}
\usepackage[margin=25mm]{geometry}
\usepackage{fontspec}
\setmainfont{Latin Modern Roman}
\usepackage{amsmath,amssymb,mathtools,unicode-math}
\setmathfont{Latin Modern Math}
\usepackage{enumitem,needspace,xurl,hyperref}
\setlength{\parindent}{0pt}
\setlength{\parskip}{5pt}
\setlength{\emergencystretch}{3em}
\newcommand{\N}{\mathbb{N}}
\newcommand{\Z}{\mathbb{Z}}
\newcommand{\Q}{\mathbb{Q}}
\newcommand{\R}{\mathbb{R}}
\newcommand{\C}{\mathbb{C}}
\newcommand{\F}{\mathbb{F}}
\newcommand{\A}{\mathbb{A}}
\newcommand{\PP}{\mathbb{P}}
\newcommand{\E}{\mathbb{E}}
\newcommand{\eps}{\varepsilon}
\newcommand{\dd}{\,\mathrm{d}}
\begin{document}
\section*{UnsolvedMath ID: 20001856}
\textbf{Source problem code:} catalogue entry
\subsection{Definitions and mathematical statement}
Thorpe's source question concerns planar corner-sharing triangles with holes of sizes five through nine and average hole size six. It calls the infinite network isostatic, and asks whether pinning every other degree-two boundary corner of a large finite piece always makes it isostatic, separately for generic and equilateral placements, and asks about joints shared by more than two bodies. The infinite-network notion, finite-piece admissibility, and pin order are not fully defined; those physical/counting descriptions alone are not a general rigidity theorem.

The aid reports a 2015 generic theorem for triangle ring networks with explicit degree, two-connectivity, boundary-cycle, and edge-cut hypotheses, but supplies neither the exact cut conditions nor a citation. For its further criterion, let \(G=(V,E)\) be a planar body-pin graph with every vertex a rigid triangle body of degree two or three, each edge a shared two-body corner, and \(B\subseteq V\) the degree-two bodies with one free corner eligible for full grounding. Assume the generically independent class and grounded sparsity characterization intended in the aid. A selected full-corner pin fixes both planar coordinates. Let \(P\subseteq B\), \(U=B\setminus P\), with \(|P|=|B|/2\), so an even boundary count is required.

For \(X\subseteq V\), write \(E(X)\) for induced edges and \(\delta(X)\) for body-graph edges crossing its cut. The reported exact generic condition is
\[
 2|P\cap X|\le|B\cap X|+|\delta(X)|,
 \quad\text{equivalently }|P\cap X|-|U\cap X|\le|\delta(X)|,
\]
for every body set, equivalently for every connected induced body set. It comes from the stated grounded count \(2|E(X)|+2|P\cap X|\le3|X|\), using degree two on \(B\) and three elsewhere. The whole-set half-boundary equality makes the total number of grounding/shared-pin constraints equal the \(3|V|\) body degrees of freedom. Counts characterize generic isostaticity only under the supplied independence/sparsity hypotheses; they do not certify a special placement. On cycle examples the aid reports nonalternating repeated PPUU patterns can work, while a proper path of three consecutive pins fails its two-edge-cut inequality.

For equilateral placement the aid reports a failure of infinitesimal isostaticity even under the same graph hypotheses: arrange six outward equilateral triangles around a regular hexagon, sharing consecutive hexagon vertices, and fix alternate free outer corners. Its grounded body-motion rigidity matrix is \(18\times18\) and reportedly has
\[
 \operatorname{rank}=15,\quad\dim\ker=3,\quad
 \dim\text{ stress space}=3.
\]
A positive-definite second-order calculation nevertheless gives local rigidity. Thus three infinitesimal mechanisms do not here imply an actual finite continuous flex. The matrix entries, exact rank certificate, and second-order form are absent; no independently verified computation is claimed. The point is the distinction between generic isostaticity, infinitesimal rigidity at a nongeneric placement, and finite/local rigidity.

The source's generalization to joints joining more than two bodies is not covered by this degree-two/degree-three, two-body-pin graph criterion. A suitable incidence/joint model and independence hypotheses would be needed; no such theorem is supplied. This partial result therefore does not settle all finite pieces of the informal source class or all equilateral networks.
\subsection{Short English statement}
The aid gives a generic cut-discrepancy test for half-boundary grounding in a specified two-body-pin class, allowing some nonalternating patterns. Its equilateral six-triangle example is infinitesimally singular but locally rigid, so generic counts cannot certify that special placement. Multi-body joints remain outside the criterion.
\subsection{Sources}
\begin{itemize}
\item Local UnsolvedMath registry, permanent ID \(20001856\), source code \texttt{catalogue entry}; \texttt{aim-geometry-notes.json}, \texttt{notes[193]}, tag \texttt{problem}.
\item AIM workshop ``Configuration spaces of linkages,'' Problem 7, attributed to Thorpe: \url{https://aimath.org/pastworkshops/linkagesproblems.pdf}. Source type PDF; statement status \texttt{exact}. Supplied verification says the original four-page PDF was inspected in that source run. Only the supplied record was used here; no PDF was accessed or compiled. The leading ``2Now'' is an extraction/page artifact.
\item The aid reports a 2015 triangle-ring generic theorem without authors, title, identifier, URL, full hypotheses, or authoritative citation. It supplies no exact grounded matrix, determinant/rank certificate, second-order form, or complete proof artifacts. Original remarks and literature fields are empty.
\item Research attempt 1: \texttt{partial\_result}, status at run \texttt{partially\_solved}, rigor label \texttt{proved}, novelty confidence medium; registry status \texttt{partially\_solved}, classification \texttt{PARTIAL-PROGRESS}. Created and updated 14 August 2026. Machine-generated aid, not peer reviewed; independently verify before citation.
\item Ulam AI / UnsolvedMath Contributors; CC BY 4.0.
\end{itemize}
\end{document}
